\section{Methodology: Bayesian Networks with Latent Time Embedding}
\label{sec:bn_lte_methodology}

\subsection{Scope of This Section}

This section describes the implemented ADNI-first Bayesian Network with Latent
Time Embedding (BN-LTE) pipeline in \texttt{BayesianNetwork-SCM}. The method
covers longitudinal target construction, train-only latent disease-time
estimation, expert-constrained dynamic Bayesian-network structure,
spline-varying structural effects, penalized estimation, held-out forecasting,
dynamic edge summaries, and comparison against mechanistic baselines such as
NDM, ESM, and SIR.

The present implementation should be interpreted as a stage-aware causal
forecasting model rather than a fully identifiable interventional causal
discovery system. Causal direction is supported by longitudinal ordering:
candidate causes are measured at baseline, while targets are future rates of
change. Biological constraints further prohibit nonphysical orientations, such
as brain atrophy causing sex or APOE genotype. Future extensions can add full
posterior MCMC over graph structures, bootstrap posterior inclusion
probabilities, external transcriptomic gating, and independent cohort
replication, but these are not required to understand the core BN-LTE model
reported here.

\subsection{Preliminaries and Notation}

The empirical unit is a longitudinal subject-pair. Let $i=1,\ldots,N$ index
eligible baseline-to-follow-up pairs, and let $u(i)$ denote the corresponding
participant. Each pair has a baseline visit time $t_{i0}$, a follow-up visit
time $t_{i1}$, and elapsed time
\begin{equation}
\Delta_i
=
\frac{t_{i1}-t_{i0}}{365.25}
>0
\end{equation}
in years. Let $X_i(t_{i0}) \in \mathbb{R}^p$ be the baseline multimodal
feature vector, let $Y_i(t_{i0}) \in \mathbb{R}^q$ be the baseline value of the
target biomarkers or regions, and let $Y_i(t_{i1}) \in \mathbb{R}^q$ be the
corresponding follow-up value. The annualized rate target for region or
biomarker $j$ is
\begin{equation}
R_{ij}
=
\frac{Y_{ij}(t_{i1}) - Y_{ij}(t_{i0})}{\Delta_i},
\qquad
R_i =
\left[
R_{i1}, \ldots, R_{iq}
\right]^T .
\end{equation}

The key latent variable is $Z_i \in [0,1]$, an estimated disease-stage
coordinate computed from baseline features only. Low values correspond to
earlier biological disease stage and high values correspond to more advanced
stage. The variable $Z_i$ is not a hidden calendar time and is not fitted using
the held-out follow-up targets; it is a train-derived, explainable ordering of
baseline disease burden.

For each target $j$, BN-LTE learns a stage-indexed directed graph
$G(z)=(V,E(z))$, where the node set $V$ contains baseline variables and future
rate targets, and $E(z)$ contains directed baseline-to-rate dependencies that
may vary with disease stage. The admissible parent set of target $j$ at stage
$z$ is
\begin{equation}
\operatorname{pa}_z(j)
=
\left\{
l : X_l(t_0) \rightarrow R_j
\text{ is active at stage } z
\right\}.
\end{equation}

\begin{center}
\begin{tabular}{p{0.22\linewidth}p{0.70\linewidth}}
Symbol & Meaning \\
\hline
$X_{il}(t_{i0})$ & Baseline value of predictor $l$ for pair $i$. \\
$Y_{ij}(t_{i0})$ & Baseline burden of target $j$. \\
$Y_{ij}(t_{i1})$ & Follow-up burden of target $j$. \\
$R_{ij}$ & Annualized future rate of target $j$. \\
$Z_i$ & Train-derived latent disease stage in $[0,1]$. \\
$G(z)$ & Disease-stage-specific directed graph. \\
$\operatorname{pa}_z(j)$ & Parents allowed to directly affect target-rate $R_j$ at stage $z$. \\
$a_j(z)$ & Stage-specific intercept or population drift term. \\
$c_j(z)$ & Stage-specific self-history effect of baseline $Y_j(t_{i0})$ on its future rate. \\
$b_{jl}(z)$ & Stage-specific direct effect of baseline predictor $X_l(t_{i0})$ on future rate $R_j$. \\
$B_k(z)$ & Cubic spline basis function evaluated at disease stage $z$. \\
$K$ & Second-difference roughness penalty matrix for spline coefficients. \\
$\Phi_j$ & Regression design matrix for target $j$. \\
\end{tabular}
\end{center}

This notation separates three temporal roles. Baseline covariates $X(t_{i0})$
are candidate causes, baseline targets $Y(t_{i0})$ provide autoregressive
disease state, and future rates $R$ are the effects to be predicted and
explained. This separation is central: BN-LTE does not infer causality from
simultaneous correlation alone.

\subsection{Algorithm 1: BN-LTE Training and Evaluation Pipeline}

\paragraph{Input.}
Raw ADNI tables, participant identifiers, visit dates, baseline multimodal
variables, regional target variables, biological graph constraints, spline
specification, candidate regularization strengths, and train/validation/test
seed.

\paragraph{Output.}
Fitted BN-LTE model, held-out follow-up predictions, dynamic edge curves,
performance metrics, and visualization artifacts.

\begin{enumerate}
\item Construct longitudinal target pairs. For each participant, order valid
visits by acquisition date, select baseline-follow-up pairs with positive
elapsed time, compute $\Delta_i$ in years, and compute
$R_{ij}=(Y_{ij}(t_{i1})-Y_{ij}(t_{i0}))/\Delta_i$.

\item Build the baseline design matrix. Assemble $X_i(t_{i0})$ from
demographics, APOE, fluid biomarkers, amyloid-PET, tau-PET, MRI-derived
neurodegeneration measures, and cognitive scores. Store baseline targets
$Y_i(t_{i0})$, follow-up targets $Y_i(t_{i1})$, and rates $R_i$. Remove
variables with insufficient training coverage or near-zero training variance.

\item Split by participant. Partition unique participants into training,
validation, and test sets so that no participant contributes pairs to more
than one split.

\item Estimate latent disease time $Z$. Select pseudotime features $M_i$ from
baseline variables only. In tau-free mode, exclude tau-PET and p-tau217
variables from $M_i$. Impute and standardize $M_i$ using training statistics,
fit a one-dimensional training SVD embedding, orient the embedding by an
interpretable disease-burden proxy, and map projected scores to $Z_i \in [0,1]$
using training percentiles.

\item Define admissible graph structure. Assign variables to biological layers:
exogenous roots, molecular pathology, imaging pathology, neurodegeneration,
and clinical outcomes. Prohibit incoming edges into root variables, prohibit
outgoing edges from clinical outcome variables, prohibit edges that violate the
baseline-to-future temporal order, and for each target $j$ rank admissible
parents by coverage, biological priority, and training association.

\item Fit stage-varying structural equations. For each target $j$, evaluate
the spline basis $B(Z_i)$ on training rows; construct $\Phi_j$ from stage
intercept terms $B_k(Z_i)$, baseline self-history terms
$Y_{ij}(t_{i0})B_k(Z_i)$, and admissible parent terms
$X_{il}(t_{i0})B_k(Z_i)$; fit penalized MAP/ridge coefficients on training
rows; select regularization by validation performance; and recover
$a_j(z)$, $c_j(z)$, and $b_{jl}(z)$ on a dense $Z$ grid.

\item Predict held-out progression. For each validation or test pair, compute
$Z_i$ using the fitted training embedding, predict future rates
$\widehat{R}_{ij}$ from the fitted structural equation, and convert rates to
follow-up burden:
\begin{equation}
\widehat{Y}_{ij}(t_{i1})
=
Y_{ij}(t_{i0}) + \Delta_i \widehat{R}_{ij}.
\end{equation}

\item Evaluate and summarize. Compute rate and follow-up-burden MAE/RMSE,
compute spatial rank agreement between empirical and predicted regional burden
maps, compare BN-LTE against NDM, ESM, SIR, persistence, and related models,
and export dynamic edge curves, causal-order summaries, and brain-region
visualizations comparing baseline, empirical follow-up, and model predictions.
\end{enumerate}

\subsection{Overview}

We model Alzheimer's disease progression as a non-stationary causal process in
which the dependency structure between biomarkers changes with disease stage.
Let subject-pair $i$ contain a baseline visit at time $t_{i0}$ and a follow-up
visit at time $t_{i1}$, with elapsed time $\Delta_i>0$ measured in years. For
each regional tau-PET target $j$, we observe baseline burden
$Y_{ij}(t_{i0})$ and follow-up burden $Y_{ij}(t_{i1})$. The learning target is
the annualized longitudinal rate
\begin{equation}
R_{ij}
=
\frac{Y_{ij}(t_{i1}) - Y_{ij}(t_{i0})}{\Delta_i}.
\label{eq:bnlte_rate}
\end{equation}

The proposed Bayesian Network with Latent Time Embedding (BN-LTE) therefore
does not attempt to orient a same-visit cross-sectional graph. Instead, every
candidate cause is measured at baseline and every outcome is a future rate of
change. This baseline-to-rate formulation gives the graph an explicit temporal
direction and avoids leakage from follow-up biomarkers into the predictor set.

For each subject-pair, BN-LTE estimates an explainable latent disease time
$Z_i \in [0,1]$ from baseline multimodal measurements. Conditional on $Z_i$,
the model defines a directed Bayesian-network structural causal model over
future rates. The key assumption is that the parent set and causal effect size
of a biomarker can be stage dependent: an edge that is weak early in disease can
strengthen later, and an early edge can decouple after a downstream pathology
becomes self-sustaining.

\subsection{Data Representation}

Each pair is represented by a baseline feature vector
\begin{equation}
X_i(t_{i0})
=
\left[
X_{i1}(t_{i0}), \ldots, X_{ip}(t_{i0})
\right]^T,
\end{equation}
containing demographic, genetic, fluid, amyloid-PET, tau-PET, MRI, and
cognitive variables. In the current ADNI implementation, root variables include
age, sex, education, and APOE $\varepsilon 4$ dosage; fluid variables include
plasma p-tau217, A$\beta$42/40, NfL, and GFAP; pathology variables include
amyloid-PET burden and regional tau-PET SUVR; neurodegeneration variables
include FreeSurfer-derived hippocampal, amygdala, and temporal cortical
measures; and clinical variables include ADAS-Cog-13, MMSE, RAVLT immediate
recall, and CDR-SB.

For tau forecasting, the response vector consists of annualized rates in a
selected temporo-parietal target set,
\begin{equation}
R_i =
\left[
R_{i1}, \ldots, R_{iq}
\right]^T,
\end{equation}
where each $R_{ij}$ is computed from the subject-specific follow-up interval
$\Delta_i$. In the model-comparison experiment, the target regions are
bilateral entorhinal, fusiform, inferior temporal, middle temporal, and inferior
parietal cortex. Subject-level train, validation, and test splits are made by
unique participant identifier so that no subject contributes pairs to more than
one split.

\subsection{Latent Time Embedding}

The latent time variable $Z_i$ is intended to represent biological disease
stage rather than chronological age. It is learned using training rows only and
is then applied unchanged to validation and test rows. This avoids using
held-out outcomes to define the disease axis.

Let $M_i \subset X_i(t_{i0})$ be the subset of baseline features permitted for
pseudotime estimation. In the default tau-free mode, all tau and p-tau217
features are excluded from $M_i$, so the latent disease coordinate used to
predict tau rates is not directly constructed from the tau targets themselves.
Features with insufficient training coverage or near-zero training variance are
excluded. Missing values are imputed with training medians, and retained
features are standardized using training means and standard deviations:
\begin{equation}
\widetilde{M}_{im}
=
\frac{M_{im}-\mu_m}{s_m}.
\end{equation}

BN-LTE then fits a one-dimensional linear embedding by taking the leading right
singular vector of the standardized training matrix:
\begin{equation}
\widetilde{M}_{\mathrm{train}}
=
U \Sigma V^T,
\qquad
w = V_{\cdot 1}.
\end{equation}
The raw latent-time score is the projection
\begin{equation}
s_i
=
\widetilde{M}_i^T w .
\end{equation}

Because singular vectors are sign ambiguous, the direction of $w$ is oriented
so that $s_i$ is positively correlated with a disease-burden proxy. The burden
proxy assigns positive sign to variables expected to increase with disease
burden, such as amyloid, NfL, GFAP, ADAS-Cog-13, and CDR-SB, and negative sign
to variables expected to decrease with disease burden, such as MMSE, RAVLT,
MRI volume, cortical thickness, and A$\beta$42/40. Scores are mapped to the
unit interval using the first and ninety-ninth training percentiles:
\begin{equation}
Z_i
=
\operatorname{clip}
\left(
\frac{s_i - Q_{0.01}(s_{\mathrm{train}})}
     {Q_{0.99}(s_{\mathrm{train}})-Q_{0.01}(s_{\mathrm{train}})},
0, 1
\right).
\label{eq:bnlte_z}
\end{equation}
This construction makes $Z_i$ explainable: each subject's position can be
decomposed into feature-level contributions $\widetilde{M}_{im}w_m$, and
population-level loadings identify which biomarkers define the disease axis.

\subsection{Dynamic Bayesian-Network Factorization}

For a fixed latent disease stage $z$, BN-LTE defines a directed acyclic
Bayesian network over future rates:
\begin{equation}
p(R_i \mid X_i(t_{i0}), Y_i(t_{i0}), Z_i=z, G_z)
=
\prod_{j=1}^{q}
p
\left(
R_{ij}
\mid
Y_{ij}(t_{i0}),
X_{i,\operatorname{pa}_z(j)}(t_{i0}),
Z_i=z
\right).
\label{eq:bnlte_factorization}
\end{equation}
Here $\operatorname{pa}_z(j)$ is the set of admissible baseline parents for
target $j$ at latent time $z$. In the current implementation, a target has a
fixed candidate parent set selected before estimation, while the corresponding
effect functions can become effectively active or inactive across $z$. Thus,
graph non-stationarity is represented by smooth stage-varying effect curves
rather than by independently estimating a separate graph at every time point.

The conditional distribution for each target is a Gaussian structural equation:
\begin{equation}
R_{ij}
=
a_j(Z_i)
+ c_j(Z_i)Y_{ij}(t_{i0})
+ \sum_{\ell \in \operatorname{pa}(j)}
b_{j\ell}(Z_i)X_{i\ell}(t_{i0})
+ \epsilon_{ij},
\qquad
\epsilon_{ij}\sim\mathcal{N}(0,\sigma_j^2).
\label{eq:bnlte_scm}
\end{equation}
The term $a_j(Z_i)$ is the stage-specific baseline rate for target $j$;
$c_j(Z_i)Y_{ij}(t_{i0})$ is a self-history term that captures autonomous
progression from existing local burden; and $b_{j\ell}(Z_i)$ is the direct
stage-dependent effect of baseline parent $\ell$ on the future rate of target
$j$. The coefficient curve $b_{j\ell}(z)$ is interpreted as a conditional
baseline-to-future-rate effect, not as a randomized interventional effect.

After estimating the rate, follow-up burden is predicted by the explicit Euler
update
\begin{equation}
\widehat{Y}_{ij}(t_{i1})
=
Y_{ij}(t_{i0}) + \Delta_i \widehat{R}_{ij}.
\label{eq:bnlte_euler}
\end{equation}

\subsection{Spline Parameterization of Non-Stationary Effects}

To permit smooth changes in effect magnitude across latent disease time, BN-LTE
represents each stage-varying function with a spline basis. Let
$B_1(z),\ldots,B_K(z)$ be cubic B-spline basis functions on $[0,1]$. Then
\begin{align}
a_j(z)
&=
\sum_{k=1}^{K}\alpha_{jk}B_k(z),
\\
c_j(z)
&=
\sum_{k=1}^{K}\kappa_{jk}B_k(z),
\\
b_{j\ell}(z)
&=
\sum_{k=1}^{K}\theta_{j\ell k}B_k(z).
\end{align}

For each target $j$, the design matrix concatenates three blocks: spline basis
columns for the stage-specific intercept, interactions between the target's
baseline value and the spline basis for self-history, and interactions between
each candidate parent and the spline basis for directed parent effects:
\begin{equation}
\Phi_{ij}
=
\left[
B(Z_i),
Y_{ij}(t_{i0})B(Z_i),
\{X_{i\ell}(t_{i0})B(Z_i)\}_{\ell \in \operatorname{pa}(j)}
\right].
\label{eq:bnlte_design}
\end{equation}
This expansion turns the non-stationary structural equation into a regularized
linear model in spline coefficients while preserving a nonlinear dependence on
disease stage.

\subsection{Biological Graph Constraints}

BN-LTE constrains candidate edge orientations using minimal disease-agnostic
biological rules. Variables are assigned to ordered layers:
\begin{equation}
\mathrm{root}
\prec
\mathrm{fluid}
\prec
\mathrm{pathology}
\prec
\mathrm{neurodegeneration}
\prec
\mathrm{clinical}.
\end{equation}
Root variables, including age, sex, education, and APOE genotype, cannot have
incoming edges. Clinical variables are treated as sinks and cannot serve as
parents. Cross-layer edges are allowed only from earlier to later layers. Within
the pathology layer, amyloid-to-tau edges are allowed, but tau-to-amyloid edges
are not. Within the fluid layer, a coarse ordering is imposed from amyloid/A
beta measures to p-tau217, GFAP, and NfL. These constraints do not hard-code a
complete Alzheimer's cascade; they only rule out biologically implausible
orientations such as a cognitive score causing genotype or future tau burden
causing baseline amyloid.

For each target, the admissible parent set is further screened by training-set
coverage and ranked by association with the target rate. A small set of
biologically motivated priority parents is prepended for tau-rate targets,
including plasma p-tau217, amyloid-PET burden, A$\beta$42/40, APOE
$\varepsilon 4$ dosage, and age. The final parent set is capped at a fixed
maximum size to control variance in the high-dimensional spline design.

\subsection{Estimation as Penalized Maximum A Posteriori Learning}

For target $j$, let $r_j$ be the vector of finite training rates and $\Phi_j$
the corresponding design matrix. BN-LTE estimates spline coefficients with
ridge-regularized least squares:
\begin{equation}
\widehat{\beta}_j
=
\arg\min_{\beta_j}
\left\{
\|r_j-\Phi_j\beta_j\|_2^2
+\lambda_j\|\beta_j\|_2^2
\right\}.
\label{eq:bnlte_ridge}
\end{equation}
Equivalently, this is the maximum a posteriori estimator under a Gaussian noise
model and an isotropic Gaussian prior on the spline coefficients:
\begin{equation}
r_j \mid \beta_j,\sigma_j^2
\sim
\mathcal{N}(\Phi_j\beta_j,\sigma_j^2 I),
\qquad
\beta_j \sim \mathcal{N}(0,\lambda_j^{-1}I).
\end{equation}
The regularization strength $\lambda_j$ is chosen by group-aware
cross-validation on the training set, where folds are grouped by subject
identifier. This keeps multiple longitudinal pairs from the same participant
inside the same fold and prevents subject-level leakage during hyperparameter
selection.

The closed-form solution is
\begin{equation}
\widehat{\beta}_j
=
(\Phi_j^T\Phi_j+\lambda_j I)^{-1}\Phi_j^T r_j,
\label{eq:bnlte_solution}
\end{equation}
after median imputation and training-set standardization of design columns.
Although the model is formulated as a Bayesian-network SCM, the current
implementation uses this penalized MAP estimator rather than full posterior
MCMC over graph structures. Posterior graph sampling is therefore a natural
extension for the next experimental phase, while the present estimator provides
a stable, leakage-controlled, interpretable approximation suitable for model
comparison and hypothesis generation.

\subsection{Edge Curves and Causal Summaries}

After fitting, the coefficient blocks are contracted with the spline basis on a
dense grid $z\in[0,1]$ to recover stage-specific effect curves:
\begin{equation}
\widehat{b}_{j\ell}(z)
=
\sum_{k=1}^{K}
\widehat{\theta}_{j\ell k}B_k(z).
\end{equation}
For comparability across variables with different units, edge curves are
reported in standardized units:
\begin{equation}
\widetilde{b}_{j\ell}(z)
=
\widehat{b}_{j\ell}(z)
\frac{\operatorname{sd}(X_\ell)}
     {\operatorname{sd}(R_j)}.
\label{eq:bnlte_standardized_edge}
\end{equation}
An edge is considered active at latent time $z$ when
$|\widetilde{b}_{j\ell}(z)|$ exceeds a pre-specified threshold. The resulting
activation profile identifies whether a candidate mechanism is early, late,
transient, or persistent across the disease continuum. Optional subject-level
bootstrapping repeats parent selection and coefficient estimation on resampled
training subjects, producing an empirical edge-stability score. These summaries
are used for mechanistic interpretation, but they should be read as
regularized conditional effect summaries rather than formal posterior inclusion
probabilities until the full graph-sampling extension is implemented.

\subsection{Evaluation Protocol}

Model evaluation is performed on held-out subjects. The primary prediction
targets are annualized tau-PET rates and follow-up tau SUVR values in the
selected target regions. Rate metrics measure whether the model captures the
speed and direction of longitudinal change, while follow-up SUVR metrics measure
absolute burden prediction:
\begin{equation}
\operatorname{MAE}_{\mathrm{rate}}
=
\frac{1}{Nq}
\sum_{i,j}
\left|
R_{ij}-\widehat{R}_{ij}
\right|,
\end{equation}
\begin{equation}
\operatorname{MAE}_{\mathrm{SUVR}}
=
\frac{1}{Nq}
\sum_{i,j}
\left|
Y_{ij}(t_{i1})-\widehat{Y}_{ij}(t_{i1})
\right|.
\end{equation}
Spatial Spearman correlation across regions is computed per subject to assess
whether the predicted spatial pattern of burden matches the empirical pattern.
Delta Spearman correlation is computed on the change vector
$Y_i(t_{i1})-Y_i(t_{i0})$ to test whether the model captures the spatial
ordering of progression beyond baseline persistence.

BN-LTE is compared against network diffusion, epidemic spreading,
susceptible-infected-recovered spreading, and baseline persistence. The
graph-spreading baselines are fit on the full 68-region tau state using the
group-average connectome, while BN-LTE is fit directly to selected regional
annualized tau rates. This comparison is intentionally conservative for
absolute follow-up SUVR, because baseline persistence already explains much of
the spatial rank structure; the more informative test for disease dynamics is
the held-out rate and delta-rank prediction.

\subsection{Interpretation}

BN-LTE should be interpreted as a stage-aware, longitudinal causal forecasting
model. The temporal design gives every edge the direction ``baseline state
predicts future rate''; the latent time embedding lets the same edge vary across
disease stage; and the biological constraints rule out implausible orientations
without hard-coding a complete disease cascade. A positive edge from baseline
amyloid burden to future tau rate, for example, means that among subjects with
similar modeled covariates and latent disease stage, higher baseline amyloid is
associated with a higher subsequent tau accumulation rate in that target
region.

The model does not by itself establish randomized interventional causality,
because ADNI is observational and unmeasured confounding remains possible.
Nevertheless, its leakage-controlled longitudinal orientation, interpretable
latent time, stage-varying effects, and explicit graph constraints make it a
useful framework for generating and testing mechanistic hypotheses about the
ordering of amyloid, tau, neurodegeneration, and cognition.
